% SPDX-License-Identifier: CC-BY-SA-4.0
% Author: Matthieu Perrin

\begingroup

\begin{exercice}[Révisions -- Minimisation d'automates] \label{exo:lexing/minimisation/revision}

  \begin{question}

  \item Donner l'automate minimal équivalent à l'automate déterministe suivant par la méthode de Moore.
    \begin{center}
      \begin{tikzpicture}[automaton, size=20mm]
          \state[initial]   (0) at (0,1) {$0$}; 
          \state[accepting] (1) at (1,1) {$1$}; 
          \state            (2) at (1,0) {$2$}; 
          \state[accepting] (3) at (2,0) {$3$}; 
          \state[accepting] (4) at (2,1) {$4$}; 
          \state            (5) at (0,0) {$5$}; 

          \path (3) edge[loop right] node       {$b$}    (3);
          \path (4) edge[loop right] node       {$b$}    (4);
          \path (5) edge[loop left]  node       {$a, b$} (5);
          \path (0) edge             node       {$a$}    (1);
          \path (0) edge             node       {$b$}    (5);
          \path (1) edge             node       {$a$}    (2);
          \path (1) edge             node       {$b$}    (4);
          \path (2) edge             node       {$b$}    (5);
          \path (4) edge             node[swap] {$a$}    (2);
          \path (3) edge[bend left]  node       {$a$}    (2);
          \path (2) edge[bend left]  node       {$a$}    (3);
      \end{tikzpicture}
    \end{center}

  \item Minimiser l'automate suivant.
    \begin{center}
    \begin{tikzpicture}[automaton, size=20mm, baseline=(1.base)]
      \state[initial, accepting] (1) at (0,1) {$1$}; 
      \state                     (2) at (1,1) {$2$}; 
      \state                     (3) at (2,1) {$3$}; 
      \state                     (4) at (0,0) {$4$}; 
      \state                     (5) at (1,0) {$5$}; 
      \state[accepting]          (6) at (2,0) {$6$}; 

      \path (1) edge[loop above] node {$a$}   (1);
      \path (4) edge[loop below] node {$a,b$} (4);
      \path (5) edge[loop below] node {$b$}   (5);
      \path (1) edge             node {$b$}   (2);
      \path (1) edge             node {$a$}   (4);
      \path (2) edge             node {$a$}   (3);
      \path (2) edge             node {$b$}   (4);
      \path (3) edge[bend left]  node {$a$}   (6);
      \path (6) edge[bend left]  node {$a$}   (3);
      \path (6) edge             node {$a$}   (5);
    \end{tikzpicture}
    \end{center}

    \item Montrer qu'un automate minimal a au plus un état stérile.
    
  \end{question}
  
\end{exercice}

\begin{correction}

  \begin{question}

  \item
    \begin{tikzpicture}[automaton, baseline=(A.base), x=20mm]
      \state[initial above] (A) at (1,0) {$A$};
      \state                (B) at (0,0) {$B$};
      \state[accepting]     (C) at (2,0) {$C$};

      \path (A) edge[bend left]  node {$a$} (C);
      \path (A) edge             node {$b$} (B);
      \path (B) edge[loop above] node {$a,b$} (B);
      \path (C) edge[bend left]  node {$a$} (A);
      \path (C) edge[loop above] node {$b$} (C);
    \end{tikzpicture}
    avec  $A = \{0,2\}$, $B = \{5\}$ et $C = \{1,3,4\}$

  \item
    \begin{tikzpicture}[automaton, baseline=(A.base), x=20mm, y=20mm]
      \state[initial, accepting] (A) at (0,1) {$A$};
      \state                     (B) at (1,1) {$B$};
      \state                     (C) at (2,1) {$C$};
      \state[accepting]          (D) at (2,0) {$D$};
      \state                     (E) at (1,0) {$E$};

      \path (A) edge[loop below] node {$a$} (A);
      \path (A) edge             node {$b$} (B);

      \path (B) edge             node {$a$} (C);
      \path (B) edge             node {$b$} (E);

      \path (C) edge[bend left]  node {$a$} (D);
      \path (C) edge             node {$b$} (E);

      \path (D) edge             node {$a$} (C);
      \path (D) edge             node {$b$} (E);

      \path (E) edge[loop left]  node {$a,b$} (E);
    \end{tikzpicture}
    
    avec $A = \{1,14\}$, $B = \{2,24\}$, $C = \{3,34,345,35\}$, $D = \{6,64\}$ et $E = \{4,45,5,\bot\}$
    
  \item Soient $p$ et $q$ deux états stériles d’un automate déterministe minimal.
    Alors $\mathcal{L}(p)=\mathcal{L}(q)=\emptyset$ (langages droits).
    Donc $p$ et $q$ sont équivalents ; par minimalité, $p=q$.
    Ainsi, un automate minimal a au plus un état stérile.

  \end{question}

\end{correction}

\endgroup
\endinput
